\section{Flow and matching}

\entry{Dinic (max flow)}{$O(n^2 m)$, unit caps $O(m\sqrt m)$ \quad mem $n + m$}{%
\code{addEdge} returns the edge id, and after \code{maxflow} the flow actually pushed
along it is \code{flowOn(id)} (only meaningful when it was added with \code{rcap = 0}).
For an undirected edge pass \code{rcap = cap}.\\[0.3ex]
\begin{ctab}{0.22}{0.73}
\code{cut()} & after \code{maxflow}: \code{c[v]} true iff \code{v} is on the
\code{s} side of a minimum cut. The cut edges are those with \code{c[u] \&\& !c[v]}.\\
\end{ctab}\\[0.3ex]
\textbf{Recipes.} Bipartite matching: source$\to$left cap 1, left$\to$right cap 1,
right$\to$sink cap 1; with unit caps this runs in $O(m\sqrt n)$, i.e. the same bound as
Hopcroft-Karp, so there is no separate template here.
Min vertex cover in a bipartite graph = max matching (K\"onig), and max independent set
= $n$ minus that.
Vertex capacities: split \code{v} into \code{v\_in -> v\_out} with that capacity.
Minimum path cover of a DAG = $n$ minus max matching on the split graph.
Project selection / max closure: source$\to$profit nodes, cost nodes$\to$sink,
prerequisites as \code{INF} edges, answer = total profit $-$ max flow.\\[0.3ex]
\code{dfs} recurses to depth $O(n)$; fine at contest sizes, but it is the thing to
rewrite if $n$ is huge.}

%LABEL Dinic
\begin{lstlisting}
struct Dinic {
    struct E { int to; ll cap; };
    int n; vector<E> e; vector<vector<int>> g;
    vector<int> lvl, it;
    Dinic(int n = 0) : n(n), g(n) {}
    int addEdge(int u, int v, ll cap, ll rcap = 0) {
        g[u].push_back((int)e.size()); e.push_back({v, cap});
        g[v].push_back((int)e.size()); e.push_back({u, rcap});
        return (int)e.size() - 2;               // edge id
    }
    ll dfs(int u, int t, ll f) {
        if (u == t || !f) return f;
        for (int& i = it[u]; i < (int)g[u].size(); i++) {
            int id = g[u][i];
            if (lvl[e[id].to] != lvl[u] + 1 || !e[id].cap)
                continue;
            if (ll d = dfs(e[id].to, t, min(f, e[id].cap))) {
                e[id].cap -= d; e[id ^ 1].cap += d;
                return d;
            }
        }
        return 0;
    }
    ll maxflow(int s, int t) {
        ll fl = 0;
        while (true) {
            lvl.assign(n, -1); lvl[s] = 0;
            vector<int> q{s};
            for (int i = 0; i < (int)q.size(); i++)
                for (int id : g[q[i]])
                    if (e[id].cap && lvl[e[id].to] < 0) {
                        lvl[e[id].to] = lvl[q[i]] + 1;
                        q.push_back(e[id].to);
                    }
            if (lvl[t] < 0) return fl;
            it.assign(n, 0);
            while (ll d = dfs(s, t, INF)) fl += d;
        }
    }
    vector<char> cut() {          // true = s side of min cut
        vector<char> c(n);
        for (int i = 0; i < n; i++) c[i] = lvl[i] >= 0;
        return c;
    }
    ll flowOn(int id) { return e[id ^ 1].cap; }
};
\end{lstlisting}

\entry{Min cost max flow}{$O(\mathrm{flow} \cdot nm)$ \quad mem $n + m$}{%
SPFA-based, so \textbf{negative edge costs are fine} (negative \emph{cycles} are not).
Returns \code{\{flow, cost\}}. Pass \code{maxf} to cap the flow, which is how you ask
for ``cheapest way to ship exactly $k$ units''.\\[0.3ex]
For \emph{min cost flow of any size} rather than min cost \emph{max} flow, stop as soon
as \code{d[t] >= 0}: further augmentation only adds cost.\\[0.3ex]
Assignment problems fit here (left$\to$right with cost, unit caps) but
\code{hungarian} is $O(n^3)$ and much faster for a dense square matrix.}

%LABEL MCMF
\begin{lstlisting}
struct MCMF {
    struct E { int to; ll cap, cost; };
    int n; vector<E> e; vector<vector<int>> g;
    MCMF(int n = 0) : n(n), g(n) {}
    void addEdge(int u, int v, ll cap, ll cost) {
        g[u].push_back((int)e.size());
        e.push_back({v, cap, cost});
        g[v].push_back((int)e.size());
        e.push_back({u, 0, -cost});
    }
    pair<ll,ll> run(int s, int t, ll maxf = INF) {
        ll fl = 0, cost = 0;
        while (fl < maxf) {
            vi d(n, INF); vector<int> pe(n, -1);
            vector<char> inq(n, 0);
            deque<int> q{s}; d[s] = 0; inq[s] = 1;
            while (!q.empty()) {              // SPFA
                int u = q.front(); q.pop_front(); inq[u] = 0;
                for (int id : g[u]) {
                    if (!e[id].cap) continue;
                    int v = e[id].to;
                    if (d[u] + e[id].cost < d[v]) {
                        d[v] = d[u] + e[id].cost; pe[v] = id;
                        if (!inq[v]) {
                            inq[v] = 1; q.push_back(v);
                        }
                    }
                }
            }
            if (d[t] >= INF) break;
            ll push = maxf - fl;
            for (int v = t; v != s; v = e[pe[v] ^ 1].to)
                push = min(push, e[pe[v]].cap);
            for (int v = t; v != s; v = e[pe[v] ^ 1].to) {
                e[pe[v]].cap -= push;
                e[pe[v] ^ 1].cap += push;
            }
            fl += push; cost += push * d[t];
        }
        return {fl, cost};
    }
};
\end{lstlisting}

\entry{Bipartite matching (Kuhn)}{$O(nm)$ \quad mem $n + m$}{%
Left side \code{[0, n1)}, right side \code{[0, n2)}, \code{addEdge(u, v)} with
\code{u} left and \code{v} right. After \code{maxMatching()}, \code{matchL[u]} /
\code{matchR[v]} give the partner or \code{-1}.\\[0.3ex]
Short and fast enough for a few thousand vertices. Above that build a unit-capacity
\code{Dinic} instead --- same graph, $O(m\sqrt n)$.\\[0.3ex]
\textbf{K\"onig, minimum vertex cover.} Let $Z$ be everything reachable by alternating
paths from the \emph{unmatched} left vertices. Then a minimum vertex cover is
$(L \setminus Z) \cup (R \cap Z)$ and a maximum independent set is its complement; both
have the expected sizes.}

%LABEL Matching
\begin{lstlisting}
struct Matching {
    int n1, n2;
    vector<vector<int>> g;
    vector<int> matchL, matchR;
    vector<char> vis;
    Matching(int n1 = 0, int n2 = 0) : n1(n1), n2(n2), g(n1),
        matchL(n1, -1), matchR(n2, -1) {}
    void addEdge(int u, int v) { g[u].push_back(v); }
    bool tryKuhn(int u) {
        for (int v : g[u]) if (!vis[v]) {
            vis[v] = 1;
            if (matchR[v] == -1 || tryKuhn(matchR[v])) {
                matchL[u] = v; matchR[v] = u;
                return true;
            }
        }
        return false;
    }
    int maxMatching() {
        int r = 0;
        for (int u = 0; u < n1; u++) {
            vis.assign(n2, 0);
            if (tryKuhn(u)) r++;
        }
        return r;
    }
};
\end{lstlisting}

\entry{Minimum vertex cover (K\"onig)}{KACTL \quad $O(n + m)$ after matching}{%
The construction above as code, KACTL's \code{cover} on our \code{Matching}. Call
\code{maxMatching()} first. Returns the cover as \{left vertices, right vertices\};
its size equals the matching. Everything \emph{not} in it is a maximum independent set
--- ``choose the most items with no conflicting pair'' when conflicts only go between
two groups.}

%DEP Matching
%LABEL cover
\begin{lstlisting}
pair<vi, vi> vertexCover(Matching& M) {
    vector<char> lf(M.n1, 1), seen(M.n2, 0);
    for (int u : M.matchR) if (u != -1) lf[u] = 0;
    vi q, L, R;
    rep(i,0,M.n1) if (lf[i]) q.push_back(i);
    while (!q.empty()) {
        int i = q.back(); q.pop_back();
        lf[i] = 1;
        for (int e : M.g[i]) if (!seen[e] && M.matchR[e]+1) {
            seen[e] = 1;
            q.push_back(M.matchR[e]);
        }
    }
    rep(i,0,M.n1) if (!lf[i]) L.push_back(i);
    rep(i,0,M.n2) if (seen[i]) R.push_back(i);
    return {L, R};
}
\end{lstlisting}

\entry{Hungarian (assignment)}{$O(n^2 m)$ \quad mem $nm$}{%
Minimum-cost \textbf{perfect} matching on a cost matrix, rows $\le$ columns: every row
gets its own column, cheapest in total. Returns \code{\{cost, asg\}} with
\code{asg[row] = column}. Input is 0-indexed; the internals are 1-indexed, which is why
the matrix is read as \code{a[i0-1][j-1]}.\\[0.3ex]
\textbf{Maximise} instead: negate every cost, then negate the answer. Forbidden pairs:
a large finite cost, not \code{INF} (\code{INF} arithmetic overflows here) --- something
like $10^{12}$ with real costs below $10^6$.\\[0.3ex]
Not the same thing as \code{Matching}: that one maximises the \emph{number} of pairs,
this one fixes the number and minimises the \emph{cost}. Costs may be negative.}

%LABEL hungarian
\begin{lstlisting}
// a = cost matrix, a.size() <= a[0].size()
pair<ll, vector<int>> hungarian(const matrix& a) {
    int n = (int)a.size(), m = (int)a[0].size();
    // 1-indexed inside; column 0 is virtual, p[j]==0 => free
    vi u(n + 1, 0), v(m + 1, 0), minv(m + 1);
    vector<int> p(m + 1, 0), way(m + 1, 0);
    for (int i = 1; i <= n; i++) {
        p[0] = i; int j0 = 0;
        fill(all(minv), INF);
        vector<char> used(m + 1, 0);
        do {
            used[j0] = 1;
            int i0 = p[j0], j1 = 0; ll delta = INF;
            for (int j = 1; j <= m; j++) if (!used[j]) {
                ll cur = a[i0-1][j-1] - u[i0] - v[j];
                if (cur < minv[j]) {
                    minv[j] = cur; way[j] = j0;
                }
                if (minv[j] < delta) {
                    delta = minv[j]; j1 = j;
                }
            }
            for (int j = 0; j <= m; j++) {
                if (used[j]) {
                    u[p[j]] += delta; v[j] -= delta;
                } else minv[j] -= delta;
            }
            j0 = j1;
        } while (p[j0] != 0);
        do {                          // flip the path back
            int j1 = way[j0]; p[j0] = p[j1]; j0 = j1;
        } while (j0);
    }
    vector<int> asg(n);
    for (int j = 1; j <= m; j++) if (p[j]) asg[p[j]-1] = j-1;
    ll cost = 0;
    for (int i = 0; i < n; i++) cost += a[i][asg[i]];
    return {cost, asg};
}
\end{lstlisting}

\entry{Eulerian path}{$O(n + m)$ \quad mem $n + m$}{%
Hierholzer, iterative, handles both directions with the \code{dir} flag. Returns the
\textbf{vertex} sequence, length \code{m+1}, using every edge exactly once --- or an
empty vector when none exists (wrong degrees, or the edges are disconnected). A closed
tour comes back with \code{path.front() == path.back()}.\\[0.3ex]
Existence: \emph{directed} needs \code{out-in} equal to 0 everywhere except possibly one
$+1$ (the start) and one $-1$ (the end); \emph{undirected} needs zero or exactly two odd
vertices. Both also need all edges in one connected component --- which the final
\code{size() != m+1} test catches, so isolated vertices do not matter.\\[0.3ex]
For the edge sequence instead, record \code{id} alongside each pop.}

%LABEL eulerPath
\begin{lstlisting}
vector<int> eulerPath(int n, const vector<pii>& es,
                      bool dir) {
    int m = (int)es.size();
    if (!m) return {};
    vector<vector<pii>> g(n);          // (to, edge id)
    vi outd(n, 0), ind(n, 0);
    for (int i = 0; i < m; i++) {
        int a = (int)es[i].first, b = (int)es[i].second;
        g[a].push_back({b, i}); outd[a]++; ind[b]++;
        if (!dir) {
            g[b].push_back({a, i}); outd[b]++; ind[a]++;
        }
    }
    int s = -1;
    if (dir) {
        int plus = 0, minus = 0;
        for (int v = 0; v < n; v++) {
            ll df = outd[v] - ind[v];
            if (df == 1) { plus++; s = v; }
            else if (df == -1) minus++;
            else if (df) return {};
        }
        if (plus > 1 || minus > 1) return {};
    } else {
        int odd = 0;
        for (int v = 0; v < n; v++)
            if (outd[v] & 1) { odd++; s = v; }
        if (odd != 0 && odd != 2) return {};
    }
    if (s < 0)
        for (int v = 0; v < n; v++)
            if (outd[v]) { s = v; break; }
    vector<int> it(n, 0), st{s}, path;
    vector<char> used(m, 0);
    while (!st.empty()) {
        int u = st.back();
        while (it[u] < (int)g[u].size()
               && used[g[u][it[u]].second]) it[u]++;
        if (it[u] == (int)g[u].size()) {
            path.push_back(u); st.pop_back();
        } else {
            auto [v, id] = g[u][it[u]++];
            used[id] = 1; st.push_back((int)v);
        }
    }
    if ((int)path.size() != m + 1) return {};   // disjoint
    reverse(all(path));
    return path;
}
\end{lstlisting}

\entry{Global min cut (Stoer--Wagner)}{KACTL \quad $O(V^3)$}{%
The cheapest set of edges whose removal disconnects an \textbf{undirected} graph, with
\emph{no} \code{s} and \code{t} given. (If they are given, that is max flow: use
\code{Dinic}.) Running Dinic for every pair would be $V^2$ flows; this is one $O(V^3)$
pass, fine for $V \le 500$. With unit weights it is the edge connectivity.\\[0.3ex]
Input is an adjacency \textbf{matrix} of total weights --- add each edge to both
\code{mat[u][v]} and \code{mat[v][u]}, summing parallel edges. Returns
\code{\{weight, vertices on one side\}}. Each phase merges the last two vertices of a
maximum-adjacency order; the cut ``last vertex alone'' is a candidate.
KACTL's \code{int}s are \code{ll} here (\code{INT\_MIN} became \code{-INF}).}

%LABEL mincut
\begin{lstlisting}
pair<ll, vi> globalMinCut(vector<vi> mat) {
    pair<ll, vi> best = {LLONG_MAX, {}};
    int n = (int)mat.size();
    vector<vi> co(n);
    rep(i,0,n) co[i] = {i};
    rep(ph,1,n) {
        vi w = mat[0];
        size_t s = 0, t = 0;
        rep(it,0,n-ph) { // O(V^2); O(E log V) with a heap
            w[t] = -INF;
            s = t, t = max_element(all(w)) - w.begin();
            rep(i,0,n) w[i] += mat[t][i];
        }
        best = min(best, {w[t] - mat[t][t], co[t]});
        co[s].insert(co[s].end(), all(co[t]));
        rep(i,0,n) mat[s][i] += mat[t][i];
        rep(i,0,n) mat[i][s] = mat[s][i];
        mat[0][t] = -INF;
    }
    return best;
}
\end{lstlisting}
